\begin{abstract}
% IEEE Access rule: one paragraph, 150-250 words, no abbreviations, footnotes, references, equations or tables.
% Numbers: headline.yaml (test 57.74, +1.03 vs DGFusion) and tasks/DRN-261019-03/analysis.md (validation missing-sensor score 57.85, +7.92).
Vehicles and robots fuse several sensors so that scene understanding survives when one sensor fails, yet the most accurate multimodal segmentation models lose much of their accuracy when a single input is missing. We investigate why. On a synthetic driving benchmark with camera, depth, event, and laser range inputs, we compare a leading fusion model that uses depth supervision with our own model that does not, under one evaluation protocol. Both depend heavily on depth: removing it costs 23 to 39 points of mean intersection over union, more at night than in daylight. When either model is also trained on synthetically degraded inputs, this dependence largely disappears, which indicates that it is learned from clean training data rather than caused by the fusion design. Building on this finding, we present a fusion model on a frozen vision foundation backbone with sensor-specific adapters, symmetric cross-sensor attention, and a high-resolution decoding path, trained with a clean pass and a degraded pass distilled from a clean teacher. The training needs no condition labels and adds no inference cost. The model reaches 57.7 mean intersection over union on the test set, one point above the best published fusion model, and 57.9 under the standard missing-sensor protocol on the validation set, eight points above the best published result. We further analyse sensor dependence across conditions, thin objects, and failure cases \todo{add one sentence on the real-sensor benchmarks once measured}.
\end{abstract}

\begin{keywords}
% IEEE Access template: three or four keywords or phrases, alphabetical order
Foundation models, missing modality, multimodal semantic segmentation, robustness, sensor fusion
\end{keywords}
